\section{Découverte~$\boldsymbol{\beta}$~: quotient stable de complexité minimale}
\label{sec:beta}

\subsection{Contexte et source}
Candidat récupéré du \textbf{Cahier~1, chapitre~XVI, page~50}. Cette page est notable car c'est la première à avoir subi une remédiation manuelle des tactiques Lean avant que la re-exécution haute génération ne réussisse.

\subsection{Eta-quotient découvert}
\begin{equation}
  f_\beta(\tau) = q^{-12/24} \cdot \eta(3\tau)^{1} \, \eta(6\tau)^{1} \, \eta(7\tau)^{-2} \, \eta(8\tau)^{4} \, \eta(9\tau)^{-6} \, \eta(10\tau)^{4} \, \eta(11\tau)^{-1} \, \eta(12\tau)^{-1}.
\end{equation}
\textbf{Énergie RAMA~:} $E = 1{,}7840$. \quad \textbf{Erreur~:} $I = 0{,}1538$.

\subsection{Calcul exact des invariants}
\begin{proposition}
$\ceff \approx 0{,}2734$, $k = 0$, $p = -5/12$.
\end{proposition}

\begin{remark}[Signification de $k=0$]
Les eta-quotients de poids~0 sont d'un intérêt particulier car ils correspondent à des fonctions algébriques sur des courbes modulaires. La célèbre fraction continue de Rogers--Ramanujan est elle-même un rapport d'eta-quotients de poids~0~\cite{andrews_berndt_I}.
\end{remark}

\subsection{Interprétation physique}
\begin{itemize}
  \item Entropie d'états BPS~: $S_{\mathrm{BPS}} = 2\pi \approx 6{,}2832$.
  \item Charge centrale effective~: $\ceff \approx 0{,}2734 > 0$.
  \item \textbf{État d'équilibre~: Stable (Unitaire)}.
\end{itemize}
